% !PRINTED none
\begin{center}
\gesperrt{Dritter Abschnitt.}

\gesperrt{Die Gruppe der Kreistheilungsk\"orper.}

\rule{3cm}{0.4pt}

\S.\,17.

\gesperrt{Die Resolventen der Kreistheilungstheorie.}
\end{center}

Von den S\"atzen \"uber \gesperrt{Abel'sche} Gruppen machen wir eine
Anwendung auf die Kreistheilungstheorie f\"ur den Fall, dass der
Grad der Einheitswurzeln nicht eine Primzahl, sondern eine
h\"ohere Potenz einer Primzahl ist. Ist $q$ eine ungerade Primzahl
und $m = q^{\varkappa}$, $\varkappa > 1$, so nehmen wir eine primitive Congruenzwurzel
 $g$ von $m$ und setzen f\"ur jede durch $q$ nicht theilbare
 Zahl $n$
\weq{1}{n \equiv g^{\nu} \pmod{m}.}
Dann ist $\nu$ der \gesperrt{Index} von $n$. Durchl\"auft $n$ die Gruppe $\mathfrak{N}$
der durch $q$ nicht theilbaren Zahlclassen nach dem Modul $m$
vom Grade
\weq{2}{c = \varphi(q^{\varkappa}) = q^{\varkappa - 1}(q - 1),}
so durchl\"auft $\nu$ ein volles Restsystem nach dem Modul $c$. Wir
nehmen nun eine primitive $m^{\text{te}}$ Einheitswurzel $r$ und eine primitive
 $c^{\text{te}}$ Einheitswurzel $\varepsilon$ und bilden die \gesperrt{Lagrange'schen}
\gesperrt{Resolventen}
\weq{3}{(\varepsilon^{\beta}, r) = \sum^{n} \varepsilon^{\beta \nu} r^n,}
worin $\beta$ ein beliebiger Exponent ist. Es handelt sich um die
Frage, wann eine solche Resolvente verschwinden kann. Ist
$\varkappa = 1$, also $m$ eine Primzahl, so verschwindet sie, wie wir im
Bd.\ I, \S.\,169 gesehen haben, f\"ur keinen Werth von $\beta$. Im allgemeinen
 Falle eines beliebigen $\varkappa$ bedeute $n'$ eine beliebige Zahl
aus $\mathfrak{N}$. Dann durchl\"auft $n n'$ zugleich mit $n$ die ganze Gruppe $\mathfrak{N}$.
Man kann also in (3) $n$ durch $n n'$ ersetzen, wenn man zugleich
%CONT>>
